The first replay corpus contains 43 requests recorded from the Astropy tasks \texttt{12907} and \texttt{13033}. Each configuration receives the same saved messages and tool definitions, sequentially with at most one active request. Reported vLLM prompt lengths range from 23,024 to 68,557 tokens. Sampling uses temperature 0.6, top-$p$ 0.95, top-$k$ 20, and presence penalty 1.0, with a 1,024-token response cap and no per-request seed. This protocol measures fixed-request serving behavior; generated tool calls are not executed. Penalties use per-path histories (Section~\ref{sec:sampling-validation}).

Native autoregressive decoding, chain MTP, and \sys{} use the same checkpoint, vLLM image, tokenizer, chat template, and common serving settings, including prefix caching, chunked prefill, a 131,072-token context limit, and one active sequence. The native arms use their native attention invocation; \sys{} uses its specialized tree-attention path. For the cross-stack comparison, we capture the complete input token sequence from vLLM's tokenizer endpoint and pass those IDs directly to SGLang's generation endpoint. All 43 sequences match the vLLM prompt lengths and the SGLang run records. SGLang EAGLE uses the same checkpoint, FlashInfer attention, and its own prefill and streaming implementation.

The disjoint corpus selects 43 of 61 captured requests from native-MTP agent executions on tasks \texttt{14096}, \texttt{14309}, \texttt{14508}, and \texttt{14995}, across five conversations. Selection allocates requests proportionally across conversations and spaces them through capture order. Its manifest is frozen before any replay; neither message bodies nor initial user messages overlap the first corpus. \sys{} and native MTP-5 replay these saved messages in three runs each; SGLang has one run. SGLang uses its chat endpoint for this corpus, so matching messages does not establish matching serialized input IDs. Input collection is distinct from the 30-attempt task-outcome study.

For vLLM and the disjoint corpus's SGLang chat endpoint, the replay client measures time from the first nonempty streamed content, reasoning, or tool-call delta to stream completion. For the first corpus's SGLang token-ID endpoint, timing begins at the first streamed increase in the completion-token count. We apply Eq.~\ref{eq:pooled-decode} to these client durations and server-reported completion counts, pooling matched token and time totals across requests and repetitions. The estimate excludes first-output latency and server startup. Subsequent stream delivery and parser behavior remain in the measured duration; the endpoint observables are not identical.
